\worksheettitle
  {Домашняя работа: ответы}
  {\modulemark{M03-H} Версия учителя}

\begin{teacherbox}[Назначение]
  Домашняя работа проверяет деление записи на классы, понимание начальных
  нулей, сборку числа и объяснение ошибок.
\end{teacherbox}

\section{Ответы}

\begin{teacherbox}[1. Разделите на классы]
  \begin{enumerate}
    \item \(3\,807\);
    \item \(92\,064\);
    \item \(705\,009\);
    \item \(6\,040\,308\);
    \item \(81\,002\,040\);
    \item \(206\,000\,015\).
  \end{enumerate}
\end{teacherbox}

\begin{teacherbox}[2 и 3. Разбор числа]
  Для \(64\,008\,205\): группы \(64\), \(008\), \(205\); количества единиц
  классов \(64\), \(8\), \(205\). Группа \(8\) заняла бы только одну позицию и
  сдвинула последующие цифры в другие разряды. Поэтому внутри числа нужна
  группа \(008\).
\end{teacherbox}

\begin{teacherbox}[4. Соберите числа]
  \begin{enumerate}
    \item \(28\,|\,004\,|\,007=28\,004\,007\);
    \item \(6\,|\,000\,|\,053=6\,000\,053\);
    \item \(104\,|\,030\,|\,009=104\,030\,009\).
  \end{enumerate}
\end{teacherbox}

\newpage
\begin{teacherbox}[5. Исправьте ошибки]
  \begin{enumerate}
    \item \(8040205=8\,|\,040\,|\,205=8\,040\,205\). Запись ошибочно
      разделили слева.
    \item \(61\,|\,002\,|\,040=61\,002\,040\). Группы тысяч и единиц не
      дополнили нулями до трёх цифр.
  \end{enumerate}
\end{teacherbox}

\begin{teacherbox}[6. Сравните]
  \(34\,006\,500<34\,060\,050\), различается класс тысяч.

  \(7\,205\,008>7\,025\,800\), различается класс тысяч.
\end{teacherbox}

\begin{teacherbox}[7. Итог]
  \begin{enumerate}
    \item \(50703009=50\,|\,703\,|\,009=50\,703\,009\);
    \item группа класса единиц \(009\), количество единиц \(9\);
    \item \(19\,|\,006\,|\,004=19\,006\,004\).
  \end{enumerate}
\end{teacherbox}

\begin{resultbox}[Критерии проверки]
  \begin{itemize}
    \item \textbf{Освоено:} не менее 15 из 17 элементов выполнены верно,
      задания 5 и 7 объяснены.
    \item \textbf{Нужна точечная доработка:} одна повторяющаяся ошибка; дать
      два аналогичных примера.
    \item \textbf{Нужен M03-R:} ученик начинает деление слева, теряет нулевые
      позиции или не выполняет обратное действие.
  \end{itemize}
\end{resultbox}
